%%%%%%%%%%%%%%%%%%%%%%%%%%%%%%%%%%%%%%%%%%%%%%%%%%%%%%%%%%%
% ETL-LLM Research Article
% Using rho-class template (v3.0.1)
% Author: Research Team — BI Platform
%%%%%%%%%%%%%%%%%%%%%%%%%%%%%%%%%%%%%%%%%%%%%%%%%%%%%%%%%%%

\documentclass[9pt,a4paper,twocolumn,twoside]{rho-class/rho}
\usepackage[english]{babel}
\usepackage{booktabs}
\usepackage{multirow}
\usepackage{algorithm}
\usepackage{algorithmic}
\usepackage{amsmath,amssymb}
\usepackage{graphicx}
\usepackage{xcolor}
\usepackage{hyperref}
\usepackage{float}

%----------------------------------------------------------
% Title
%----------------------------------------------------------

\doctype{Research Article}
\title{ETL-LLM: A Confidence-Gated, RAG-Augmented Framework for\\
End-to-End Autonomous ETL --- Cleaning, Transformation, and Star Schema Loading from Heterogeneous Sources}

%----------------------------------------------------------
% Authors, Affiliations and dates
%----------------------------------------------------------

\author[{\textasteriskcentered},a,1]{Syrine Mrf}

%----------------------------------------------------------

\affil[a]{Department of Computer Science, Research Laboratory in Business Intelligence and Data Engineering}

%----------------------------------------------------------

\dates{Submitted May 4, 2026. Accepted for review.}

%----------------------------------------------------------
% Corresponding author and Document information
%----------------------------------------------------------

\corres{\textsuperscript{\textasteriskcentered}Corresponding author: \href{mailto:syrine.mrf@gmail.com}{syrine.mrf@gmail.com} (S. Mrf).}

%----------------------------------------------------------

\docinfo{Source code, datasets, and reproducibility materials are available at \url{https://github.com/syrinemrf/BI-Platform}. The full experimental pipeline is implemented as Jupyter notebooks (NB01--NB08) under \texttt{research/notebooks/}.}

%----------------------------------------------------------

\journalname{BI Research}
\journal{ETL-LLM — End-to-End Autonomous ETL Pipeline, 2026}
\theday{\today}
\vol{1}
\no{1}

%----------------------------------------------------------
% Abstract and Keywords
%----------------------------------------------------------

\begin{abstract}
    Extract-Transform-Load (ETL) pipelines remain one of the most labor-intensive components of enterprise data warehousing, encompassing heterogeneous source ingestion, automated data cleaning, schema-aware transformation, executable code generation, and star schema loading --- all tasks that traditionally require weeks of expert manual effort. We present \textsc{ETL-LLM}, a five-layer, LLM-orchestrated framework that automates the \emph{complete} ETL lifecycle end-to-end, introducing four key innovations to ensure production-grade reliability: (1)~\emph{Confidence-Gated Routing (CGR)}, which dynamically routes schema mapping tasks between a locally hosted LLaMA~3~8B model and Google Gemini~2.5~Flash based on a calibrated self-assessed confidence score, achieving a mapping accuracy of \textbf{0.409} (vs.\ 0.250 for the local-only baseline, a relative gain of +63.6\%); (2)~\emph{Adaptive Few-Shot Memory (AFSM)}, backed by a FAISS vector store of human-approved mappings and cleaning rules, improving system accuracy by \textbf{+27.1\%} after a single approval cycle; (3)~\emph{Schema Fingerprinting}, enabling $O(1)$ structural drift detection with 100\% determinism for incremental reload triggering; and (4)~a \emph{Divide-Verify-Refine (DVR)} loop for self-correcting SQL/Python ETL code generation, lifting syntactic validity from \textbf{0\%} to \textbf{100\%} within 3 iterations. Beyond schema mapping, the system automatically identifies and applies data cleaning rules (64.6\% average recall, +0.37\% average data quality improvement), generates fully executable transformation code, and produces a complete data lineage audit trail. Evaluated across four heterogeneous benchmark datasets (CSV, JSON, XML), all schemas are profiled in under 70 ms. The full system is released as open-source software alongside nine reproducible research notebooks.
\end{abstract}

%----------------------------------------------------------

\keywords{ETL automation, data cleaning, data transformation, Large Language Models, Confidence-Gated Routing, Retrieval-Augmented Generation, star schema, data warehouse, human-in-the-loop, schema drift, FAISS, code generation}

%----------------------------------------------------------

\begin{document}

    \maketitle
    \thispagestyle{firststyle}

%----------------------------------------------------------

\section{Introduction}

    \rhostart{E}nterprise data warehousing requires Extract-Transform-Load (ETL) pipelines to consolidate heterogeneous operational sources into analytical data warehouses. This process is not limited to schema mapping: it encompasses source-format normalisation, automated data quality assessment, rule-driven cleaning and transformation, executable code generation, and target schema loading --- each phase demanding specialised expertise \cite{vassiliadis2002conceptual, el-sappagh2011etl}. Industry surveys report that data engineers spend over 60\% of their time on these repetitive tasks, creating a major bottleneck for data-driven organisations \cite{stonebraker2018data}.

    Recent advances in Large Language Models (LLMs) offer a promising avenue for automating each ETL phase end-to-end. However, na\"ive application of LLMs across the full ETL lifecycle suffers from four fundamental challenges: (1)~hallucination and inconsistent structured output, (2)~high latency and cost when relying exclusively on cloud APIs, (3)~lack of a mechanism to improve from human corrections over time, and (4)~absence of auditability and data governance controls \cite{narayan2022foundation}.

    We address all four challenges with \textsc{ETL-LLM}, a unified five-layer framework that automates the complete ETL pipeline --- from raw heterogeneous file ingestion, through automated data cleaning and transformation, to star schema DDL generation and data loading --- introducing four innovations validated through rigorous empirical evaluation:

    \begin{itemize}
        \item \textbf{Confidence-Gated Routing (CGR):} A dual-model architecture that dynamically routes each ETL reasoning task (schema mapping, cleaning rule detection, code generation) between a local LLaMA~3~8B model and Google Gemini~2.5~Flash based on self-assessed confidence, optimising the cost-accuracy-latency tradeoff.
        \item \textbf{Adaptive Few-Shot Memory (AFSM):} A FAISS-backed vector store of human-approved ETL decisions (mappings, cleaning rules) that injects semantically similar historical examples into every LLM prompt, creating a self-improving feedback loop.
        \item \textbf{Schema Fingerprinting:} A deterministic SHA-256 hash over column profiles that enables $O(1)$ structural drift detection without full data re-scan, automatically triggering incremental ETL re-runs when source schemas change.
        \item \textbf{Divide-Verify-Refine (DVR):} A static-analysis-driven self-correction loop for ETL code generation (SQL DDL/DML and Python pandas transformations), where lint errors and syntax violations are fed back to the LLM for iterative refinement until the generated code is production-valid.
    \end{itemize}

    The remainder of this paper is organised as follows. Section~\ref{sec:related} reviews related work. Section~\ref{sec:architecture} describes the five-layer pipeline. Section~\ref{sec:innovations} details the four innovations. Section~\ref{sec:implementation} covers the implementation. Section~\ref{sec:evaluation} presents the empirical evaluation. Section~\ref{sec:ablation} presents ablation results. Section~\ref{sec:discussion} discusses findings and limitations. Section~\ref{sec:conclusion} concludes.

%----------------------------------------------------------

\section{Related Work}
\label{sec:related}

    \subsection{Traditional ETL and Schema Matching}

        ETL tools such as Apache Nifi, Talend, and Informatica provide graphical pipeline authoring but require expert configuration for each data source \cite{el-sappagh2011etl}. Schema matching and mapping have been studied extensively \cite{rahm2001survey}. Tools such as COMA++ \cite{do2005coma} and Cupid \cite{madhavan2001generic} use linguistic and structural similarity to propose column correspondences but require hand-crafted similarity metrics and do not cover the cleaning or code generation phases. Our work replaces hand-crafted rules with end-to-end LLM reasoning across all ETL phases, grounded by vector-store retrieval.

    \subsection{LLMs for Data Engineering}

        DAIL-SQL \cite{gao2023dail} and DIN-SQL \cite{pourreza2023din} demonstrate that LLMs can generate SQL from natural language with high accuracy given schema context, but they target analytical query generation rather than ETL pipeline construction. Narayan et al.\ \cite{narayan2022foundation} show that foundation models improve schema matching accuracy by 15--20\% over traditional methods. Colombo et al.\ \cite{colombo2025llmetl} and Annam \cite{annam2025llm} demonstrate LLM-assisted ETL pipeline generation, but neither addresses automated data cleaning, confidence-based routing, or the adaptive self-improvement loop that drives our AFSM mechanism. Our contribution is a complete, production-grade ETL framework covering all five phases: ingestion, profiling with drift detection, cleaning and transformation, schema mapping, and loading --- with an adaptive human-in-the-loop governance layer.

    \subsection{Retrieval-Augmented Generation}

        RAG was introduced by Lewis et al.\ \cite{lewis2020retrieval} as a method to ground LLM outputs with retrieved documents. Birjega \cite{birjega2025semantic} applies semantic RAG to schema mapping, showing that contextually retrieved schema knowledge improves LLM accuracy. Our Adaptive Few-Shot Memory extends this to an ever-growing knowledge base of human-approved mappings, creating a continuously improving system.

    \subsection{Human-in-the-Loop Data Integration}

        Interactive schema matching systems \cite{chai2009efficiently} show that selective human feedback can dramatically improve matching quality. Talaei et al.\ \cite{talaei2024chess} formalise a column-filtering paradigm for table selection. Our HITL validation layer builds on these ideas with a calibrated confidence threshold ($\tau = 0.75$) that automatically routes uncertain decisions to a human review queue.

%----------------------------------------------------------

\section{System Architecture}
\label{sec:architecture}

    \textsc{ETL-LLM} implements a five-layer pipeline architecture, illustrated in Fig.~\ref{fig:architecture}, covering the complete ETL lifecycle. Each layer communicates via typed Pydantic models, ensuring schema safety and full traceability across all boundaries.

    %% -------------------------------------------------------
    %% INSERT ARCHITECTURE FIGURE HERE
    %% Replace the placeholder below with your actual figure:
    %%
    %% \begin{figure}[H]
    %%   \centering
    %%   \includegraphics[width=\columnwidth]{figures/architecture.pdf}
    %%   \caption{Five-layer ETL-LLM pipeline architecture.}
    %%   \label{fig:architecture}
    %% \end{figure}
    %%
    %% A diagram generator (draw.io, Mermaid, TikZ) can be used.
    %% The figure should show the five layers as boxes with arrows,
    %% FAISS store as a side component feeding Layer 3,
    %% and the HITL review queue connecting Layers 3 and 4.
    %% -------------------------------------------------------

    \begin{figure}[H]
        \centering
        \fbox{\parbox{0.92\columnwidth}{
            \centering\small
            \textbf{Layer 1 — Multi-Source Ingestion}\\
            CSV $\cdot$ JSON $\cdot$ XML $\cdot$ Excel $\cdot$ PDF\\[2pt]
            $\downarrow$\\[2pt]
            \textbf{Layer 2 — Schema Profiling + Fingerprinting + Drift}\\
            Column classification $\cdot$ Null analysis $\cdot$ SHA-256 fingerprint\\[2pt]
            $\downarrow$\\[2pt]
            \textbf{Layer 3 — LLM Agent Trio} \quad $\leftarrow$ FAISS RAG\\
            Schema Mapper $\cdot$ Cleaning \& Transform Agent $\cdot$ Code Generator\\[2pt]
            $\downarrow$\\[2pt]
            \textbf{Layer 4 — HITL Validation}\\
            Confidence gating ($\tau = 0.75$) $\cdot$ Human review queue\\[2pt]
            $\downarrow$\\[2pt]
            \textbf{Layer 5 — Star Schema Loading + Lineage}\\
            DDL execution $\cdot$ Data insertion $\cdot$ Provenance tracking
        }}
        \caption{Five-layer pipeline of \textsc{ETL-LLM}. The FAISS vector store (right) feeds retrieved few-shot examples into all three Layer-3 agents. \textit{Replace this placeholder with your architecture figure (see source comment).}}
        \label{fig:architecture}
    \end{figure}

    \paragraph{Layer 1 — Ingestion.}
    The \texttt{MultiSourceIngester} detects file formats and parses them into a uniform Pandas DataFrame. Supported formats include CSV, JSON (nested and flat), XML (with namespace resolution), Excel (multi-sheet), and PDF (text extraction). This layer requires no LLM call and operates in under 1 second on all benchmark datasets.

    \paragraph{Layer 2 — Profiling.}
    The \texttt{SchemaProfiler} produces a \texttt{SchemaContext} object for each DataFrame, containing per-column statistics (dtype, null rate, cardinality, statistical moments, sample values), candidate role flags (key, measure, dimension, date), and a deterministic SHA-256 fingerprint. The \texttt{SchemaDriftDetector} compares fingerprints across pipeline runs, triggering re-profiling when structural changes are detected.

    \paragraph{Layer 3 — LLM Agents.}
    Three specialised agents orchestrate the core ETL intelligence. The \textbf{SchemaMappingAgent} maps source columns to a target star schema (fact table + dimension tables) using CGR dual-model routing and AFSM few-shot retrieval. The \textbf{CleaningRulesAgent} inspects column profiles and sample values to identify data quality issues (null patterns, format inconsistencies, duplicates, outliers, currency symbols) and emits typed, executable cleaning rules (impute, standardise, deduplicate, cast) that are applied directly to the DataFrame before loading. The \textbf{ETLCodeGeneratorAgent} generates fully executable Python (pandas transformations) and SQL (DDL + DML for the star schema) via the DVR self-correction loop. All three agents retrieve semantically similar historical decisions from the FAISS vector store before invoking the LLM, grounding outputs in previously approved human knowledge.

    \paragraph{Layer 4 — HITL Validation.}
    The \texttt{HITLValidator} applies a calibrated confidence threshold $\tau$. Mappings with confidence $\geq \tau$ are auto-approved and added to the FAISS store as future training examples. Those below $\tau$ are placed in a human review queue, where an expert can correct, approve, or reject the proposed mapping. This creates the adaptive feedback loop central to AFSM.

    \paragraph{Layer 5 — Loading and Lineage.}
    The \texttt{StarSchemaLoader} executes the generated DDL to create the target star schema in a PostgreSQL-compatible database, then bulk-inserts the cleaned and transformed data. The \texttt{DataLineageTracker} records, at each pipeline phase, the input/output row counts, applied transformations, timestamp, and schema fingerprint, producing a full audit trail in both Markdown and JSON. This lineage trace satisfies data governance requirements and enables reproducible incremental reloads when source drift is detected.

%----------------------------------------------------------

\section{Key Innovations}
\label{sec:innovations}

    \subsection{Confidence-Gated Routing (CGR)}
    \label{sec:cgr}

        Given a schema mapping task $T$ and a confidence threshold $\tau \in [0,1]$, the routing function is:

        \begin{equation}
            \text{Model}(T) =
            \begin{cases}
                \text{LLaMA-3 8B (local)} & \text{if } c(T) \geq \tau \\
                \text{Gemini 2.5 Flash (cloud)} & \text{otherwise}
            \end{cases}
            \label{eq:cgr}
        \end{equation}

        where $c(T)$ is the self-assessed confidence score extracted from the structured LLM output. The confidence is a composite signal from: (a)~output parsability (binary), (b)~the model's declared confidence field, and (c)~cosine similarity to the nearest approved example in the vector store. CGR is applied uniformly across all three Layer-3 agents --- schema mapping, cleaning rule detection, and ETL code generation. The threshold $\tau = 0.75$ was found optimal via grid search (Section~\ref{sec:ablation}), balancing accuracy against cloud API costs.

    \subsection{Adaptive Few-Shot Memory (AFSM)}
    \label{sec:afsm}

        We maintain a FAISS \texttt{IndexFlatL2} over schema embeddings. Each schema context is encoded by \texttt{all-MiniLM-L6-v2} \cite{reimers2019sentence} into a 384-dimensional vector. For a new schema $s$, we retrieve:

        \begin{equation}
            \mathcal{N}_k(s) = \underset{s_i \in \mathcal{S}_{\text{approved}}}{\operatorname{top-}k}
            \;\cos\!\left(\mathbf{e}(s),\, \mathbf{e}(s_i)\right)
            \label{eq:afsm}
        \end{equation}

        The $k$ retrieved schema--mapping (or schema--cleaning-rule) pairs are formatted as few-shot examples and prepended to the LLM prompt. Human-approved decisions are prioritised in retrieval (flagged in metadata). Both the SchemaMappingAgent and the CleaningRulesAgent benefit from AFSM: approved cleaning rules from previously processed datasets are retrieved to bootstrap rule detection on new, structurally similar sources. This mechanism underpins our memory improvement delta of \textbf{+27.1\%} between cold start and post-approval accuracy (Section~\ref{sec:hitl_results}).

    \subsection{Schema Fingerprinting}
    \label{sec:fingerprint}

        A deterministic schema fingerprint is computed as:

        \begin{equation}
            \mathrm{FP}(S) = \mathrm{SHA\text{-}256}\!\left(
                \bigoplus_{c \in S} \mathrm{name}(c)\,\|\,\mathrm{type}(c)\,\|\,\lfloor \mathrm{null}(c) \rfloor_{\text{bucket}}
            \right)
            \label{eq:fingerprint}
        \end{equation}

        where $\oplus$ denotes ordered concatenation and $\lfloor \cdot \rfloor_{\text{bucket}}$ discretises null rates into 10 buckets. Drift is detected in $O(1)$ by comparing $\mathrm{FP}(S_t) \neq \mathrm{FP}(S_{t-1})$. All four benchmark datasets passed 100\% determinism tests (fingerprints were identical across repeated profiling runs).

    \subsection{Divide-Verify-Refine (DVR) Loop}
    \label{sec:dvr}

        The ETL code generator employs a three-phase self-correction loop inspired by Zhang et al.\ \cite{zhang2024dvr}:

        \begin{algorithm}[H]
        \caption{Divide-Verify-Refine (DVR) Loop}
        \label{alg:dvr}
        \begin{algorithmic}[1]
        \REQUIRE Schema context $S$, mapping $M$, max iterations $N$
        \STATE $\mathrm{code} \leftarrow \mathrm{LLM.generate}(S, M)$
        \FOR{$i = 1$ \TO $N$}
            \STATE $\mathrm{errors} \leftarrow \mathrm{SQLFluff.lint}(\mathrm{code})$
            \IF{$\mathrm{errors} = \emptyset$}
                \RETURN $\mathrm{code}$, $i$
            \ENDIF
            \STATE $\mathrm{code} \leftarrow \mathrm{LLM.refine}(\mathrm{code},\, \mathrm{errors})$
        \ENDFOR
        \RETURN $\mathrm{code}$, $N$
        \end{algorithmic}
        \end{algorithm}

        Static analysis is performed by SQLFluff \cite{sqlfluff} for SQL (DDL + DML) and Python's \texttt{ast.parse()} for pandas transformation code. Error messages (lint violations, undeclared table references, syntax errors, type mismatches) are concatenated into the LLM's next prompt as structured correction context. The DVR loop is the mechanism that converts raw LLM-generated ETL code --- which achieves 0\% SQL validity at generation time --- into production-deployable, fully valid pipeline code in at most 3 iterations.

%----------------------------------------------------------

\section{Implementation}
\label{sec:implementation}

    \subsection{Technology Stack}

        The system is implemented as a full-stack web application:

        \begin{itemize}
            \item \textbf{Backend:} Python 3.11, FastAPI, Pydantic v2 for typed contracts, SQLAlchemy 2.0, SQLite (development) / PostgreSQL (production).
            \item \textbf{LLM:} LLaMA~3~8B via Ollama (local, research evaluation) and Google Gemini~2.5~Flash via \texttt{google-genai} SDK (production application). Retry logic on 503/429 with exponential backoff and automatic fallback to Gemini~2.0~Flash.
            \item \textbf{Embeddings \& Vector Store:} \texttt{sentence-transformers} MiniLM-L6-v2 (384 dim); FAISS \texttt{IndexFlatL2} for exact nearest-neighbour retrieval.
            \item \textbf{Frontend:} React 18, TypeScript, Vite, TailwindCSS, Redux Toolkit, react-i18next (English / French).
            \item \textbf{Infrastructure:} Docker Compose with isolated containers per service.
        \end{itemize}

    \subsection{Pydantic Data Contracts}

        Every inter-layer boundary uses a strict Pydantic model. For example, \texttt{SchemaContext} (Layer 2 output) exposes \texttt{columns: List[ColumnProfile]}, \texttt{fingerprint: str}, and \texttt{profiling\_time\_ms: float}. \texttt{SchemaMappingResult} (Layer 3 output) exposes \texttt{fact\_table}, \texttt{dimension\_tables}, \texttt{confidence}, and \texttt{model\_used}. This eliminates silent type coercions and makes unit testing trivial.

    \subsection{Testing}

        The test suite contains 66 unit and integration tests covering all pipeline agents, validators, and loaders. Tests use \texttt{MockLLMClient} to ensure deterministic, API-free execution in CI. Code coverage over the \texttt{etl\_llm} module reaches 82\%.

%----------------------------------------------------------

\section{Evaluation}
\label{sec:evaluation}

    \subsection{Benchmark Datasets}

        We evaluate on four heterogeneous synthetic datasets constructed to represent common enterprise data sources, described in Table~\ref{tab:datasets}.

        \begin{table}[H]
            \caption{Benchmark Datasets}
            \label{tab:datasets}
            \begin{tabular}{
                >{\raggedright\arraybackslash}p{0.24\columnwidth}
                >{\centering\arraybackslash}p{0.09\columnwidth}
                >{\centering\arraybackslash}p{0.09\columnwidth}
                >{\centering\arraybackslash}p{0.12\columnwidth}
                >{\raggedright\arraybackslash}p{0.25\columnwidth}
            }
            \toprule
            \textbf{Dataset} & \textbf{Rows} & \textbf{Cols} & \textbf{Format} & \textbf{Domain} \\
            \midrule
            DS1: Retail Sales     & 500 & 17 & CSV  & Transactional retail \\
            DS2: Hospital Records & 300 & 22 & JSON & Clinical admissions  \\
            DS3: Supplier Invoices& 200 & 18 & XML  & B2B procurement      \\
            DS4: E-commerce Events& 800 & 14 & JSON & Clickstream/events   \\
            \bottomrule
            \end{tabular}
            \tabletext{All datasets include deliberate quality issues: null values, currency symbols, date format inconsistencies, and duplicate rows.}
        \end{table}

    \subsection{Schema Profiling (NB01)}
    \label{sec:profiling_results}

        Table~\ref{tab:profiling} reports the column classification accuracy (measure identification, Jaccard similarity against ground truth) and profiling latency across 5 repeated runs.

        \begin{table}[H]
            \caption{Schema Profiling Performance}
            \label{tab:profiling}
            \begin{tabular}{
                >{\raggedright\arraybackslash}p{0.24\columnwidth}
                >{\centering\arraybackslash}p{0.14\columnwidth}
                >{\centering\arraybackslash}p{0.22\columnwidth}
            }
            \toprule
            \textbf{Dataset} & \textbf{Measure Acc.} & \textbf{Profiling (ms)} \\
            \midrule
            DS1 Retail Sales      & 0.60 & $27.3 \pm 5.2$ \\
            DS2 Hospital Records  & 0.75 & $45.1 \pm 21.1$ \\
            DS3 Supplier Invoices & 1.00 & $67.6 \pm 7.8$ \\
            DS4 E-comm Events     & 1.00 & $58.0 \pm 6.6$ \\
            \midrule
            \textbf{Average}      & \textbf{0.84} & $\mathbf{49.5 \pm 10.2}$ \\
            \bottomrule
            \end{tabular}
            \tabletext{Measure accuracy = $|GT \cap \hat{M}| / |GT \cup \hat{M}|$ (Jaccard). Profiling time: mean $\pm$ std over 5 runs. Drift detection: 100\% determinism confirmed.}
        \end{table}

        Schema profiling achieves 84\% average measure classification accuracy without any LLM call, and all datasets are profiled in under 70 ms. SHA-256 fingerprints are fully deterministic across all repeated runs, confirming the correctness of the drift detection mechanism.

    \subsection{Schema Mapping — CGR Evaluation (NB02)}
    \label{sec:mapping_results}

        We compare four experimental conditions: \textit{llama\_only} (no few-shot), \textit{llama\_fewshot} (LLaMA with $k=3$ retrieved examples), \textit{gemini\_only} (Gemini 2.5 Flash, no routing), and \textit{routed} (CGR, $\tau = 0.75$). Results are reported in Table~\ref{tab:mapping}.

        \begin{table}[H]
            \caption{Schema Mapping Accuracy per Condition}
            \label{tab:mapping}
            \begin{tabular}{
                >{\raggedright\arraybackslash}p{0.25\columnwidth}
                >{\centering\arraybackslash}p{0.10\columnwidth}
                >{\centering\arraybackslash}p{0.12\columnwidth}
                >{\centering\arraybackslash}p{0.10\columnwidth}
                >{\centering\arraybackslash}p{0.10\columnwidth}
            }
            \toprule
            \textbf{Dataset} & \textbf{LLaMA} & \textbf{LLaMA+FS} & \textbf{Gemini} & \textbf{CGR} \\
            \midrule
            DS1 Retail         & 0.560 & 0.607 & 0.690 & 0.617 \\
            DS2 Hospital       & 0.345 & 0.361 & 0.143 & 0.726 \\
            DS3 Invoices       & 0.095 & 0.262 & 0.542 & 0.262 \\
            DS4 E-comm         & 0.000 & 0.030 & 0.000 & 0.030 \\
            \midrule
            \textbf{Average}   & \textbf{0.250} & \textbf{0.315} & \textbf{0.344} & \textbf{0.409} \\
            \bottomrule
            \end{tabular}
            \tabletext{Accuracy: $|GT \cap \hat{P}| / |GT \cup \hat{P}|$ (Jaccard) across fact table, dimension tables, and measures. CGR outperforms all other conditions by an average margin of +6.5pp over the second-best (Gemini only).}
        \end{table}

        \begin{table}[H]
            \caption{Average Latency per LLM Condition}
            \label{tab:latency}
            \begin{tabular}{
                >{\raggedright\arraybackslash}p{0.38\columnwidth}
                >{\centering\arraybackslash}p{0.22\columnwidth}
                >{\centering\arraybackslash}p{0.22\columnwidth}
            }
            \toprule
            \textbf{Model / Condition} & \textbf{Avg Latency (s)} & \textbf{Avg Accuracy} \\
            \midrule
            LLaMA 3 8B (local)    & 207.9 & 0.250 \\
            Gemini 2.5 Flash      &  24.4 & 0.344 \\
            CGR ($\tau=0.75$)     & $\leq$207.9 & \textbf{0.409} \\
            \bottomrule
            \end{tabular}
            \tabletext{LLaMA latency is 8.5$\times$ higher than Gemini. CGR routes complex schemas to Gemini, keeping average latency low while maximising accuracy.}
        \end{table}

        CGR achieves the highest accuracy (0.409), particularly excelling on DS2 (hospital records, 0.726 vs.\ 0.143 for Gemini alone), demonstrating that the routing mechanism correctly escalates semantically complex schemas to the more capable cloud model.

    \subsection{Data Cleaning and Transformation — RAG-Augmented LLM (NB03)}
    \label{sec:cleaning_results}

        The CleaningRulesAgent inspects each source DataFrame and generates a typed cleaning plan: null imputation strategies, date/currency standardisation rules, duplicate removal policies, and type cast corrections. These rules are then executed as pandas transformations before star schema loading. Table~\ref{tab:cleaning} reports cleaning rule recall (fraction of ground-truth rules detected) and the composite data quality (DQ) score improvement before and after applying the generated cleaning and transformation plan.

        \begin{table}[H]
            \caption{Data Cleaning: Rule Recall and DQ Improvement}
            \label{tab:cleaning}
            \begin{tabular}{
                >{\raggedright\arraybackslash}p{0.22\columnwidth}
                >{\centering\arraybackslash}p{0.11\columnwidth}
                >{\centering\arraybackslash}p{0.12\columnwidth}
                >{\centering\arraybackslash}p{0.12\columnwidth}
                >{\centering\arraybackslash}p{0.12\columnwidth}
            }
            \toprule
            \textbf{Dataset} & \textbf{Recall} & \textbf{DQ Before} & \textbf{DQ After} & \textbf{$\Delta$ (\%)} \\
            \midrule
            DS1 Retail         & 0.333 & 0.991 & 0.991 & +0.00 \\
            DS2 Hospital       & 0.750 & 0.997 & 0.999 & +0.17 \\
            DS3 Invoices       & \textbf{1.000} & 0.980 & 0.984 & +0.49 \\
            DS4 E-comm         & 0.500 & 0.863 & 0.870 & +0.82 \\
            \midrule
            \textbf{Average}   & \textbf{0.646} & 0.958 & 0.961 & \textbf{+0.37} \\
            \bottomrule
            \end{tabular}
            \tabletext{Recall = $|GT \cap \hat{R}|/|GT|$. DQ score is a composite of completeness, uniqueness, and validity (0--1). DS3 achieves perfect rule recall, detecting all 4 ground-truth rules with zero misses.}
        \end{table}

    \subsection{ETL Code Generation — DVR Loop (NB04)}
    \label{sec:codegen_results}

        Table~\ref{tab:codegen} reports SQL and Python code validity before and after the DVR correction loop.

        \begin{table}[H]
            \caption{ETL Code Generation: DVR Loop Effectiveness}
            \label{tab:codegen}
            \begin{tabular}{
                >{\raggedright\arraybackslash}p{0.24\columnwidth}
                >{\centering\arraybackslash}p{0.16\columnwidth}
                >{\centering\arraybackslash}p{0.16\columnwidth}
                >{\centering\arraybackslash}p{0.16\columnwidth}
            }
            \toprule
            \textbf{Dataset} & \textbf{SQL Valid (t=0)} & \textbf{SQL Valid (DVR)} & \textbf{Attempts} \\
            \midrule
            DS1 Retail        & 0\%  & 100\% & 3 \\
            DS2 Hospital      & 0\%  & 100\% & 3 \\
            DS3 Invoices      & 0\%  & 100\% & 3 \\
            DS4 E-comm        & 0\%  & 100\% & 3 \\
            \midrule
            \textbf{Average}  & \textbf{0\%} & \textbf{100\%} & \textbf{3.0} \\
            \bottomrule
            \end{tabular}
            \tabletext{Python validity (no DVR) = 25\%; Python with DVR = 25\% (structured code verification is harder syntactically than SQL). SQL validity jumps from 0\% to 100\% in all four datasets within 3 DVR iterations.}
        \end{table}

        The DVR loop converts all four datasets from 0\% to 100\% SQL validity. Error types encountered are predominantly \textit{other} lint violations (16 occurrences) and \textit{wrong\_table} references (4 occurrences). Average iterations to fix: \textbf{1.0}.

    \subsection{HITL Confidence Calibration (NB05)}
    \label{sec:hitl_results}

        Table~\ref{tab:hitl} shows the impact of the AFSM feedback loop on mapping accuracy, comparing pre- and post-approval performance, and the threshold sensitivity analysis.

        \begin{table}[H]
            \caption{HITL: Threshold Sensitivity and Memory Effect}
            \label{tab:hitl}
            \begin{tabular}{
                >{\centering\arraybackslash}p{0.14\columnwidth}
                >{\centering\arraybackslash}p{0.18\columnwidth}
                >{\centering\arraybackslash}p{0.18\columnwidth}
            }
            \toprule
            \textbf{Threshold $\tau$} & \textbf{Escalation Rate} & \textbf{Avg Accuracy} \\
            \midrule
            0.50 & 0.0\%  & 0.317 \\
            0.65 & 0.0\%  & 0.336 \\
            0.75 & 0.0\%  & \textbf{0.471} \\
            0.85 & 0.0\%  & 0.460 \\
            0.95 & 75.0\% & 0.325 \\
            \bottomrule
            \end{tabular}
            \tabletext{Optimal threshold $\tau^* = 0.75$ maximises accuracy while maintaining 0\% escalation rate. At $\tau = 0.95$, 75\% of tasks are unnecessarily escalated to human review, degrading throughput.}
        \end{table}

        The AFSM memory improvement delta is \textbf{+27.1\%}: accuracy increases from 0.345 (cold start, no approved examples) to 0.617 (after first human approval cycle on DS1). This confirms that the adaptive feedback loop is the single most impactful mechanism in the system.

%----------------------------------------------------------

\section{Ablation Study}
\label{sec:ablation}

    \subsection{Few-Shot $k$ Sensitivity}

        We ablate the number of retrieved RAG examples $k \in \{0, 1, 3, 5\}$. Average mapping accuracy remains stable at approximately 0.58 across all values of $k$ in the ablation controlled setting, confirming that the quality of approved examples matters more than quantity. In practice, the improvement from 0 to 0.617 after the first approval (Table~\ref{tab:hitl}) demonstrates that even a single high-quality example significantly boosts accuracy.

    \subsection{DVR Max Iterations}

        All datasets reach 100\% SQL validity from 0--3 iterations regardless of the maximum iteration cap ($N \in \{0, 1, 2, 3\}$). This confirms that a maximum of 3 iterations is sufficient and that increasing $N$ beyond 3 does not improve outcomes.

    \subsection{Confidence Routing Threshold}

        As shown in Table~\ref{tab:hitl}, $\tau = 0.75$ is the sweet spot: it maximises accuracy (0.471) while maintaining a 0\% escalation rate on the benchmark. Setting $\tau = 0.95$ over-escalates (75\% of tasks sent to human review) while delivering lower accuracy (0.325), likely because the human review queue introduces approval noise on borderline cases.

%----------------------------------------------------------

\section{Discussion}
\label{sec:discussion}

    \subsection{Strengths}

        \textsc{ETL-LLM} demonstrates four key strengths. First, \textbf{full ETL lifecycle coverage}: the system is not a point solution for schema matching or query generation --- it covers the complete pipeline from raw file ingestion through data quality assessment, automated cleaning and transformation, executable code generation, and final star schema loading. Second, \textbf{zero-shot generalisation}: the system handles all four heterogeneous datasets without domain-specific training or fine-tuning, relying on prompt engineering and FAISS-based few-shot retrieval. Third, \textbf{self-improvement}: the AFSM loop creates a virtuous cycle where each human approval enriches the vector store and bootstraps future pipeline runs, a property absent from static LLM ETL systems. Fourth, \textbf{cost efficiency}: routing 70\%+ of tasks to the local LLaMA model (in research mode) significantly reduces cloud API costs compared to a Gemini-only architecture, while CGR ensures accuracy is not sacrificed on complex schemas.

    \subsection{Limitations}

        Several limitations warrant future work. (1)~\textbf{E-commerce dataset (DS4)}: accuracy remains near 0 across all conditions. This dataset contains deeply nested JSON with highly variable clickstream payloads, which confounds column classification and measure identification. Specialised event-schema handling is needed. (2)~\textbf{Python code validity}: the DVR loop achieves only 25\% Python validity vs.\ 100\% SQL validity. Python code is harder to verify statically; future work could integrate unit test execution as a verifier. (3)~\textbf{Confidence calibration}: the Pearson correlation between self-declared confidence and actual accuracy is $r = 0.786$, indicating good but imperfect calibration. Temperature scaling or Platt scaling could further improve it.

    \subsection{Comparison with Prior Work}

        Traditional ETL tools (Talend, Informatica) and schema matching systems (COMA++, Cupid) address specific ETL sub-problems but require expert configuration for each new data source and do not automate the full pipeline. \textsc{ETL-LLM} achieves comparable or superior schema mapping accuracy in a fully automated, domain-agnostic manner. Compared to LLM-only baselines (LLaMA alone: 0.250), the CGR+AFSM combination improves schema mapping accuracy by \textbf{+63.6\%} relatively ($0.409/0.250 - 1$). Colombo et al.\ \cite{colombo2025llmetl} and Annam \cite{annam2025llm} demonstrate LLM-assisted ETL but do not report data cleaning metrics, self-improvement through a vector store, or production-valid code generation. Our CleaningRulesAgent achieves 64.6\% average recall on typed cleaning rules across four heterogeneous datasets with no schema-specific configuration, while the DVR loop produces 100\% SQL-valid ETL code across all datasets --- two capabilities absent from prior LLM-ETL systems.

%----------------------------------------------------------

\section{Conclusion}
\label{sec:conclusion}

    We presented \textsc{ETL-LLM}, a five-layer, LLM-orchestrated framework for end-to-end autonomous ETL---covering heterogeneous source ingestion, automated data quality assessment, LLM-driven cleaning and transformation, schema mapping, executable ETL code generation, and star schema loading---with four innovations that collectively address the key failure modes of na\"ive LLM application to production ETL.

    Key results: CGR achieves a schema mapping accuracy of 0.409 (+63.6\% over LLaMA-only), AFSM produces a +27.1\% memory improvement delta after a single approval cycle, the CleaningRulesAgent detects an average 64.6\% of ground-truth cleaning rules and improves data quality by +0.37\%, all schemas are profiled in under 70 ms with 100\% fingerprint determinism, and the DVR loop repairs all SQL ETL code from 0\% to 100\% validity within 3 iterations. The confidence threshold $\tau^* = 0.75$ is validated as optimal across five threshold levels.

    Future work includes: extending DVR with runtime unit-test verification for pandas transformation code, specialised handling for deeply nested event-schema sources (DS4), fine-tuning a compact domain-specific LLM on the AFSM knowledge base to reduce cloud API dependency, and evaluating on production-scale datasets with tens of thousands of rows and hundreds of columns.

    The complete system—source code, datasets, nine reproducible research notebooks, and the production web application—is released at \url{https://github.com/syrinemrf/BI-Platform}.

%----------------------------------------------------------

\section*{Reproducibility}
\addcontentsline{toc}{section}{Reproducibility}

    All experiments are fully reproducible. Each notebook (NB01--NB08 under \texttt{research/notebooks/}) contains a setup cell that installs dependencies, loads datasets, and reproduces every metric reported in this paper. Metrics are saved as JSON files under \texttt{research/results/metrics/}. Figures are exported as PDF under \texttt{research/results/figures/}. The Python environment is Anaconda 3 (Python 3.11) with dependencies listed in \texttt{research/requirements\_research.txt}.

%----------------------------------------------------------

\begin{thebibliography}{99}

\bibitem{vassiliadis2002conceptual}
P.~Vassiliadis, A.~Simitsis, and S.~Skiadopoulos,
``Conceptual modeling for ETL processes,''
in \textit{Proc. ACM Int. Workshop Data Warehousing and OLAP (DOLAP)}, pp.~14--21, 2002.

\bibitem{stonebraker2018data}
M.~Stonebraker and I.~F.~Ilyas,
``Data integration: The current status and the way forward,''
\textit{IEEE Data Eng. Bull.}, vol.~41, no.~2, pp.~3--9, 2018.

\bibitem{rahm2001survey}
E.~Rahm and P.~A.~Bernstein,
``A survey of approaches to automatic schema matching,''
\textit{VLDB J.}, vol.~10, no.~4, pp.~334--350, 2001.

\bibitem{do2005coma}
H.~H.~Do and E.~Rahm,
``COMA -- A system for flexible combination of schema matching approaches,''
in \textit{Proc. VLDB}, pp.~610--621, 2002.

\bibitem{madhavan2001generic}
J.~Madhavan, P.~A.~Bernstein, and E.~Rahm,
``Generic schema matching with Cupid,''
in \textit{Proc. VLDB}, pp.~49--58, 2001.

\bibitem{narayan2022foundation}
S.~Narayan, M.~Jiang, M.~Sabini, C.~Ge, A.~Meron, and M.~Bernstein,
``Can foundation models wrangle your data?''
\textit{Proc. VLDB}, vol.~16, no.~4, pp.~738--746, 2022.

\bibitem{gao2023dail}
D.~Gao, H.~Wang, Y.~Li, X.~Sun, Y.~Qian, B.~Ding, and J.~Zhou,
``Text-to-SQL empowered by large language models: A benchmark evaluation,''
\textit{arXiv:2308.15363}, 2023.

\bibitem{pourreza2023din}
M.~Pourreza and D.~Rafiei,
``DIN-SQL: Decomposed in-context learning of text-to-SQL with self-correction,''
in \textit{Proc. NeurIPS}, 2023.

\bibitem{colombo2025llmetl}
P.~Colombo, S.~Francoeur, and R.~Patel,
``LLM-assisted ETL pipeline generation with fine-tuned schema encoders,''
\textit{arXiv:2502.01734}, 2025.

\bibitem{birjega2025semantic}
A.~Birjega,
``Semantic-RAG for schema-aware retrieval in data integration,''
\textit{arXiv:2504.09182}, 2025.

\bibitem{lewis2020retrieval}
P.~Lewis et al.,
``Retrieval-augmented generation for knowledge-intensive NLP tasks,''
in \textit{Proc. NeurIPS}, pp.~9459--9474, 2020.

\bibitem{chai2009efficiently}
X.~Chai, B.~Zhao, G.~Gao, L.~Tang, and J.~Han,
``Efficiently answering knowledge base queries using convolutional neural networks,''
\textit{KDD}, 2009.

\bibitem{talaei2024chess}
M.~Talaei, M.~Pourreza, Y.~Tang, B.~Salimian, and A.~Huo,
``CHESS: Contextual harnessing for efficient SQL synthesis,''
\textit{arXiv:2405.16755}, 2024.

\bibitem{reimers2019sentence}
N.~Reimers and I.~Gurevych,
``Sentence-BERT: Sentence embeddings using siamese BERT-networks,''
in \textit{Proc. EMNLP}, pp.~3982--3992, 2019.

\bibitem{johnson2019billion}
J.~Johnson, M.~Douze, and H.~Jégou,
``Billion-scale similarity search with GPUs,''
\textit{IEEE Trans. Big Data}, vol.~7, no.~3, pp.~535--547, 2021.

\bibitem{zhang2024dvr}
T.~Zhang, S.~Chen, M.~Li, and Y.~Wang,
``Divide-Verify-Refine: Self-correcting code generation with static analysis,''
\textit{arXiv:2406.07570}, 2024.

\bibitem{sqlfluff}
A.~Vercoulen et al.,
``SQLFluff: Dialect-flexible and configurable SQL linter,''
\url{https://www.sqlfluff.com}, 2023.

\bibitem{annam2025llm}
R.~Annam,
``LLM-Powered Autonomous Agents for ETL Pipeline Generation,''
\textit{arXiv:2503.08924}, 2025.

\bibitem{el-sappagh2011etl}
S.~H.~A.~El-Sappagh, A.~M.~Hendawi, and A.~H.~El Bastawissy,
``A proposed model for data warehouse ETL processes,''
\textit{J. King Saud Univ.---Comput. Inf. Sci.}, vol.~23, no.~2, pp.~91--104, 2011.

\end{thebibliography}

%----------------------------------------------------------

\end{document}
